\par\SourceLeaf{492}%
This coming summer he will first make the astronomical observations at the northern tip of Jutland. Then he will repeat last year’s measurements at Lauenburg with new instruments and determine the azimuths, and will measure a baseline several miles long near Hamburg. My presence is desired only for part of the work at Lauenburg and for the measurement of the baseline. My presence is not desired for the operations at the northern tip of Jutland. In asking for my presence during these delicate operations, the Danish government seems to me to have taken the view that joint deliberation would be of help in this undertaking, which is being carried out solely for the benefit of science. By the nature of the matter, this could be achieved only imperfectly or even not at all by correspondence alone. Secondly, since the whole degree survey is to be made known in full to the scholarly world in a work of its own, the presence of a witness to some of the most important operations will help strengthen the authenticity of the survey. Thus in 1798 France invited all states that were not then at war with it to send qualified astronomers to the baseline measurement near Paris. Denmark, Switzerland, Holland, Spain and other states did so. These two considerations apply in the first instance only to the Danish degree surveys taken by themselves. On the assumption that a further continuation through Hanover could be ordered sooner or later, it also seems particularly important that an astronomer from our side should be able to give due weight to the confidence especially merited by the operations to which such a continuation would have to be directly connected.

In view of these considerations I am very glad to comply with the request, which is flattering to me, and I will ask the Danish astronomer to notify me as early as possible of when those operations will take place, so that I can arrange my other affairs accordingly. This will probably cause no difficulty with my lectures, since those operations will not begin until late summer. I shall indeed regret in all likelihood being unable to observe here the great annular eclipse of the Sun on September 7, unique of its kind. I shall console myself, however, with the thought that Prof. \textsc{Harding} can still observe it here and that it will also appear annular in Holstein, where it will be easy to obtain good means for observing this remarkable phenomenon, even if they are not quite equal to those available here. — — —

Göttingen, February 27, 1820.
\begin{flushright}C. F. \textsc{Gauss}.\end{flushright}

\par\bigskip\Needspace{8\baselineskip}
\pdfbookmark[0]{Memorandum on the Hanoverian degree survey.}{pm-hannoversche-gradmessung}
\begin{center}\large MEMORANDUM ON THE HANOVERIAN DEGREE SURVEY.\end{center}
\par
— — — In my earlier report I already explained why several principal instruments are indispensable for this undertaking. As long as it was still uncertain even where they could be obtained, the operations themselves naturally could not begin, apart from some preliminary astronomical work at the observatory here, which I began last summer. In the meantime I have placed orders for the principal instruments required, in particular a larger theodolite of the highest possible perfection and a so-called universal instrument, from the foremost instrument maker now living, von \textsc{Reichenbach} in Munich. He has both accepted the order and promised to deliver these instruments at the beginning of May and in July of this year, respectively.
